% ============================================================
%  第二卷 第三章 电磁场中的电荷（§15–§25 + 习题）
%  底本：高教社中文版第8版；OCR 错误已修正，脚注文本待脚注代理补录
%  已修正的 OCR/原书错误：
%   1. §16 电荷 "$e_\cdot$"→$e$；"引人"→"引入"；OCR 丢失的粗体矢量统一 \boldsymbol
%   2. §17 (17.8)(17.9) 变换箭头 OCR 丢失，按原书补 "→"
%   3. §18 "可将 $\cdot\partial f/\partial x^k$" 中多余 "\cdot" 删除（照 (18.1) 正文应为 ∂f/∂x^k）
%   4. §20 "代人"→"代入"；"$\mathcal{E}_0^{~~}$"→$\mathcal{E}_0$
%   5. §21 "$\omega=ecH/\delta$"→$\omega=ecH/\mathcal{E}$（OCR δ 为花体 𝓔 之误）；
%      "$p\imath$"→$p_l$；OCR 段落断裂"…决定.粒子/的速度是常数"合并；
%      习题3 "$\hbar=$次项"→"二次项"、"$\bar\kappa\ \zeta(t)$"→$\boldsymbol{\zeta}(t)$、
%      "$\dot\pi$ 生"→"产生"、"(H.Alven,1940)"→(H. Alfvén, 1940)、纵向速度句语序按原书复原
%   6. §24 (24.2) E_z 分母 $\sqrt{1-V^2/c}$→$\sqrt{1-V^2/c^2}$
%   7. §25 (25.1)(25.2) 右端 OCR 乱码（"𝒦𝒱 dH/dt"等）→inv；
%      (25.6) OCR 丢失虚单位 i，照原书 p.73 补为
%      $F_y=F'_y\cosh\varphi-\mathrm{i}F'_z\sinh\varphi=F'_y\cos\mathrm{i}\varphi-F'_z\sin\mathrm{i}\varphi$，
%      $F_z=F'_z\cos\mathrm{i}\varphi+F'_y\sin\mathrm{i}\varphi$
%  遗留问题：
%   a. 脚注已转录（2026-10-10 脚注代理自原书扫描页补录，全章共10处）。
%      其中 §21 两处脚注①②未见于 脚注页码对照.txt，已直接自原书 p.75
%      （p-075.png）页脚补录；§18 脚注①在原书 p.70 页脚（对照表标 p070，
%      实际标记位于该页上半部 §18 末段）；
%   b. §16 起拉格朗日量 (16.4) 用花体 $\mathcal{L}$、(16.9) 起为斜体 $L$，照原书保留；
%   c. §21 "（像称圆轨道的中心那样）"译文照 OCR/原书保留。
% ============================================================
\chapter{电磁场中的电荷}

\section{相对论中的基本粒子}

粒子与粒子间的相互作用，可以借助力场的概念来描述.这就是说，我们不说一个粒子作用于另一个粒子，而说粒子在它自己周围建立起场；这个场内的任何其他粒子都受一定的力作用.在经典力学中，场仅仅是描述粒子相互作用这一物理现象的一种方法.在相对论中，因为相互作用是以有限速度传播的，情况就根本改变了.在某一时刻，作用在一个粒子上的力并不由其他粒子在该时刻的位置决定.粒子中的一个改变了位置，仅在经过某一段时间以后方能影响到别的粒子.这就是说，场本身具有物理的真实性.我们不可以说，彼此间有距离的粒子直接地相互作用.在每一时刻的相互作用仅能发生在空间中紧密相邻的各点之间（接触作用）.所以，我们应该说，一个粒子同场相互作用，而后场与另一个粒子相互作用.

我们将研究两种形态的场：引力场与电磁场.关于引力场，我们留到第十章至第十四章再研究，在其余各章内，我们只讨论电磁场.

在考虑粒子同电磁场的相互作用以前，我们要对相对论力学中“粒子”的概念做一些评论.

在经典力学中，人们可以引入刚体的概念，所谓刚体，就是在任何条件下都不能变形的物体.在相对论中，刚体似乎应该相应地理解为这样一些物体，在它们处于静止的参考系中其所有尺寸都保持不变.然而，不难看出，相对论使得一般情况下刚体不可能存在.

例如，我们考虑一个绕自身轴转动的圆盘，并假设它是刚体.固联于该盘的参考系显然不是惯性系.但是，对于圆盘的每个无限小单元，可以引入一个惯性参考系，其中该单元在某一给定时刻处于静止；对于圆盘上有不同速度的单元，这些惯性系显然是不同的.现在我们来考察沿盘某一半径分布的一系列线元.因为圆盘是刚体，所以每一段的长度（在与该段相应的惯性系中）与圆盘静止时该段的长度相同.因为每一段都垂直于自己的速度，因而在这种情况下就没有洛伦兹收缩，所以，一个静止的观察者，当圆盘的半径从他旁边扫过时，所量出的线段长度，与圆盘静止时所量出的一样.因此，静止观察者量得的半径总长（等于组成它的各段之和），同圆盘静止时量得的一样.另一方面，在给定时刻，圆盘圆周上从静止观察者旁边经过的每一单元的长度都要发生洛伦兹收缩，所以整个圆周的长度（即静止观察者所测出的各线段之和）将小于静止圆盘的周长.于是我们得出结论，由于圆盘的转动，（静止观察者测得的）圆周与半径之比必须改变，而不再等于 $2\pi$.这一结论的荒谬性表明，实际上圆盘不可能是刚体，它在转动时必然发生了某种复杂的形变，这种形变与其组成物质的弹性有关.

我们还可以用另一种方法来证明刚体是不可能存在的.假设某一外力作用于某一固体的某一点上，使这个物体运动.如果这个物体是刚体，那么，它的所有各点都必须与受外力作用的点一起运动，否则物体就要变形了.然而根据相对论，这是不可能的.因为力的作用是以有限速度从其作用点传到另外的点，因而，所有的点不可能同时开始运动.

从这些讨论中我们得出有关基本粒子的一些结论.基本粒子是这样的粒子，我们认为可以通过给定其作为整体的三个坐标和三个速度分量，就能完全决定它们的力学状态.显然，如果基本粒子具有有限的尺寸，即具有空间广延，那么它就不可能变形，因为变形的概念同物体各个部分独立运动的可能性联系着.但是我们刚才已经看到，相对论指出刚体是不可能存在的.

因此我们得到一个重要的结论：在经典的（非量子的）相对论力学中，我们不能赋予那些被看做基本的粒子有限的尺寸.换句话说，在经典理论的框架内，必须把基本粒子当做几何点看待\footnote{量子力学在这种情况下做出了根本的改变，但相对论在这里再次使引入任何非点的相互作用变得极其困难.}.

\section{场的四维势}

一个在给定电磁场中运动粒子的作用量由两部分组成：一项是自由粒子的作用量 (8.1)，另一项描述粒子与场的相互作用.后者必定包括表征粒子的量和表征场的量.

事实表明\footnote{下面的论断在一定程度上可以看成是实验数据的推论.粒子在电磁场中作用量的形式不能只基于一般的考虑（例如像相对论不变性的要求）来确定.后者会允许（16.1）式中出现 $\int A\,\mathrm{d}s$ 形式的量，式中 $A$ 是一个标量函数.
\par 为了避免任何误解，我们再次指出，我们现在考虑的是经典（而非量子）理论，因而不包括与粒子自旋有关的效应.}，粒子同电磁场相互作用的性质由一个量所决定.这个量称为粒子的电荷 $e$.电荷可以为正，也可以为负（也可以为零）.场的性质由一个四维矢量 $A_i$——四维势表征，其分量是坐标和时间的函数.这些量借助式子
\begin{equation*}
  -\frac{e}{c}\int_a^b A_i\,\mathrm{d}x^i
\end{equation*}
出现在作用量里，式中函数 $A_i$ 取值于粒子世界线上的点.因子 $1/c$ 是为方便引入的.应当指出，只要我们还没有把电荷或势同已知的量联系起来的公式，测量这些新量的单位就可以任意选取\footnote{关于这些单位的建立，参见 §27.}.

因此，电磁场中电荷的作用量将有如下的形式：
\begin{equation}
  S=\int_a^b\left(-mc\,\mathrm{d}s-\frac{e}{c}A_i\,\mathrm{d}x^i\right).
\end{equation}

四维矢量 $A^i$ 的三个空间分量构成一个三维空间矢量 $\boldsymbol{A}$，称为场的矢势，时间分量称为标势；我们记之为 $A^0=\varphi$.因此，
\begin{equation}
  A^i=(\varphi,\boldsymbol{A}).
\end{equation}

所以作用量积分可以写为下式：
\begin{equation*}
  S=\int_a^b\left(-mc\,\mathrm{d}s+\frac{e}{c}\boldsymbol{A}\cdot\mathrm{d}\boldsymbol{r}-e\varphi\,\mathrm{d}t\right).
\end{equation*}

引入 $\boldsymbol{v}=\mathrm{d}\boldsymbol{r}/\mathrm{d}t$，并变为对 $t$ 积分：
\begin{equation}
  S=\int_{t_1}^{t_2}\left(-mc^2\sqrt{1-\frac{v^2}{c^2}}
  +\frac{e}{c}\boldsymbol{A}\cdot\boldsymbol{v}-e\varphi\right)\mathrm{d}t.
\end{equation}

被积函数正是一个电荷在电磁场中的拉格朗日量：
\begin{equation}
  \mathcal{L}=-mc^2\sqrt{1-\frac{v^2}{c^2}}
  +\frac{e}{c}\boldsymbol{A}\cdot\boldsymbol{v}-e\varphi.
\end{equation}

这个函数与自由粒子的拉格朗日量 (8.2) 相差之项是
$(e/c)\boldsymbol{A}\cdot\boldsymbol{v}-e\varphi$，该项描述电荷与场的相互作用.

导数 $\partial\mathcal{L}/\partial\boldsymbol{v}$ 是粒子的广义动量；我们用 $\boldsymbol{P}$ 来代表它.微分可得
\begin{equation}
  \boldsymbol{P}=\frac{m\boldsymbol{v}}{\sqrt{1-\dfrac{v^2}{c^2}}}
  +\frac{e}{c}\boldsymbol{A}
  =\boldsymbol{p}+\frac{e}{c}\boldsymbol{A}.
\end{equation}

这里我们用 $\boldsymbol{p}$ 代表这个粒子的普通动量，以后我们简称之为动量.

由拉格朗日量，按照熟知的通式
\begin{equation*}
  \mathcal{H}=\boldsymbol{v}\cdot\frac{\partial\mathcal{L}}{\partial\boldsymbol{v}}-\mathcal{L},
\end{equation*}
我们可以求出粒子在场内的哈密顿量.将（16.4）代入，我们得到
\begin{equation}
  \mathcal{H}=\frac{mc^2}{\sqrt{1-\dfrac{v^2}{c^2}}}+e\varphi.
\end{equation}

但是，哈密顿量不应该用速度来表示，而应该用粒子的广义动量来表示.

从 (16.5) 和 (16.6) 显见，$\mathcal{H}-e\varphi$ 与
$\boldsymbol{P}-(e/c)\boldsymbol{A}$ 之间的关系同场不存在时
$\mathcal{H}$ 和 $\boldsymbol{p}$ 之间的关系一样，即
\begin{equation}
  \left(\frac{\mathcal{H}-e\varphi}{c}\right)^2=m^2c^2
  +\left(\boldsymbol{P}-\frac{e}{c}\boldsymbol{A}\right)^2,
\end{equation}
或者
\begin{equation}
  \mathcal{H}=\sqrt{m^2c^4+c^2\left(\boldsymbol{P}-\frac{e}{c}\boldsymbol{A}\right)^2}+e\varphi.
\end{equation}

对于低速情况，即是说在经典力学中，拉格朗日量（16.4）化为
\begin{equation}
  L=\frac{mv^2}{2}+\frac{e}{c}\boldsymbol{A}\cdot\boldsymbol{v}-e\varphi.
\end{equation}

在这种近似下，
\begin{equation*}
  \boldsymbol{p}=m\boldsymbol{v}
  =\boldsymbol{P}-\frac{e}{c}\boldsymbol{A},
\end{equation*}
并且我们可求得哈密顿量的表达式：
\begin{equation}
  \mathcal{H}=\frac{1}{2m}\left(\boldsymbol{P}-\frac{e}{c}\boldsymbol{A}\right)^2+e\varphi.
\end{equation}

最后，我们来写电磁场中粒子的哈密顿—雅可比方程.在哈密顿量中，用
$\partial S/\partial\boldsymbol{r}$ 代替 $\boldsymbol{P}$，用 $-\partial S/\partial t$ 代替
$\mathcal{H}$，就可得到这个方程.因此，由 (16.7)，我们得到
\begin{equation}
  \left(\nabla S-\frac{e}{c}\boldsymbol{A}\right)^2
  -\frac{1}{c^2}\left(\frac{\partial S}{\partial t}+e\varphi\right)^2
  +m^2c^2=0.
\end{equation}

\section{场中电荷的运动方程}

位于场内的电荷不只受到场的作用力，而是也反过来对场起作用，改变场.但是，假如电荷 $e$ 不大，电荷对于场的作用就可以略去不计.在这种情况下，当我们考虑电荷在一给定的场内运动时，我们可以假设场本身与电荷的坐标或速度无关.要能够把电荷认为是在上述意义上的小，它所应该满足的确切条件将在以后讲述（见§75）.以下我们假设这个条件已被满足.

我们现在要找出电荷在给定电磁场内的运动方程.这些方程可以通过对作用量进行变分得到.因而，运动方程就是拉格朗日方程
\begin{equation}
  \frac{\mathrm{d}}{\mathrm{d}t}\frac{\partial L}{\partial\boldsymbol{v}}
  =\frac{\partial L}{\partial\boldsymbol{r}},
\end{equation}
其中 $L$ 由公式（16.4）确定.

导数 $\partial L/\partial\boldsymbol{v}$ 是粒子的广义动量 (16.5).其次，我们可写出
\begin{equation*}
  \frac{\partial L}{\partial\boldsymbol{r}}\equiv\nabla L
  =\frac{e}{c}\operatorname{grad}\left(\boldsymbol{A}\cdot\boldsymbol{v}\right)
  -e\operatorname{grad}\varphi.
\end{equation*}

但是由已知的矢量分析的公式可得
\begin{equation*}
  \operatorname{grad}\left(\boldsymbol{a}\cdot\boldsymbol{b}\right)
  =\left(\boldsymbol{a}\cdot\nabla\right)\boldsymbol{b}
  +\left(\boldsymbol{b}\cdot\nabla\right)\boldsymbol{a}
  +\boldsymbol{b}\times\operatorname{rot}\boldsymbol{a}
  +\boldsymbol{a}\times\operatorname{rot}\boldsymbol{b},
\end{equation*}
其中 $\boldsymbol{a}$ 及 $\boldsymbol{b}$ 是两个任意矢量.应用这个公式到
$\boldsymbol{A}\cdot\boldsymbol{v}$，并记住，对 $\boldsymbol{r}$ 微分时，
$\boldsymbol{v}$ 是常数，如是，则求得
\begin{equation*}
  \frac{\partial L}{\partial\boldsymbol{r}}
  =\frac{e}{c}\left(\boldsymbol{v}\cdot\nabla\right)\boldsymbol{A}
  +\frac{e}{c}\boldsymbol{v}\times\operatorname{rot}\boldsymbol{A}
  -e\operatorname{grad}\varphi.
\end{equation*}

因此，拉格朗日方程具有如下的形式：
\begin{equation*}
  \frac{\mathrm{d}}{\mathrm{d}t}\left(\boldsymbol{p}+\frac{e}{c}\boldsymbol{A}\right)
  =\frac{e}{c}\left(\boldsymbol{v}\cdot\nabla\right)\boldsymbol{A}
  +\frac{e}{c}\boldsymbol{v}\times\operatorname{rot}\boldsymbol{A}
  -e\operatorname{grad}\varphi.
\end{equation*}

{\emergencystretch=1.5em
但是全微分 $(\mathrm{d}\boldsymbol{A}/\mathrm{d}t)\mathrm{d}t$ 包含两部分：矢势在空间的某一给定点因时间变化而发生的变化
$(\partial\boldsymbol{A}/\partial t)\mathrm{d}t$，以及由空间的一点移动一段距离
$\mathrm{d}\boldsymbol{r}$ 至另一点所发生的变化.第二部分是
$(\mathrm{d}\boldsymbol{r}\cdot\nabla)\boldsymbol{A}$.因此导数
$\mathrm{d}\boldsymbol{A}/\mathrm{d}t$ 可以写成以下的形式：
\par}
\begin{equation*}
  \frac{\mathrm{d}\boldsymbol{A}}{\mathrm{d}t}
  =\frac{\partial\boldsymbol{A}}{\partial t}
  +\left(\boldsymbol{v}\cdot\nabla\right)\boldsymbol{A}.
\end{equation*}

将此式代入前面的方程，得到
\begin{equation}
  \frac{\mathrm{d}\boldsymbol{p}}{\mathrm{d}t}
  =-\frac{e}{c}\frac{\partial\boldsymbol{A}}{\partial t}
  -e\operatorname{grad}\varphi
  +\frac{e}{c}\boldsymbol{v}\times\operatorname{rot}\boldsymbol{A}.
\end{equation}

这就是一个粒子在电磁场内的运动方程.等号左边是粒子的动量对时间的导数.所以 (17.2) 式等号右边的式子就是作用于电磁场内的粒子上的力.

我们看出，这个力包含两部分.第一部分（(17.2) 式右边的第一、第二两项）与粒子的速度无关.第二部分（第三项）与速度有关，它与速度成正比，而且垂直于速度.

作用于单位电荷上的第一种类型的力，称为电场强度；我们用 $\boldsymbol{E}$ 来代表它.于是，按定义，
\begin{equation}
  \boldsymbol{E}=-\frac{1}{c}\frac{\partial\boldsymbol{A}}{\partial t}
  -\operatorname{grad}\varphi.
\end{equation}

作用于单位电荷上的第二种类型的力中，速度的因子，严格说来是 $v/c$ 的因子，称为磁场强度.我们用 $\boldsymbol{H}$ 来代表它.于是，按定义，
\begin{equation}
  \boldsymbol{H}=\operatorname{rot}\boldsymbol{A}.
\end{equation}

假如在一电磁场中，$\boldsymbol{E}\neq0$，但 $\boldsymbol{H}=0$，我们就说它是电场；假如
$\boldsymbol{E}=0$ 但 $\boldsymbol{H}\neq0$，我们就说它是磁场.在一般情况下，电磁场是电场与磁场的叠加.

我们指出，$\boldsymbol{E}$ 是极矢量，而 $\boldsymbol{H}$ 是轴矢量.

一个电荷在电磁场中的运动方程现在可以写成
\begin{equation}
  \frac{\mathrm{d}\boldsymbol{p}}{\mathrm{d}t}
  =e\boldsymbol{E}+\frac{e}{c}\boldsymbol{v}\times\boldsymbol{H}.
\end{equation}

等号右边的式子称为洛伦兹力.其第一部分（电场作用于电荷上的力）与电荷速度无关，并沿着 $\boldsymbol{E}$ 的方向.第二部分（磁场作用于电荷上的力）与电荷的速度成正比，而其方向则既垂直于速度又垂直于磁场 $\boldsymbol{H}$.

对于比光速小很多的速度，动量 $\boldsymbol{p}$ 近似地等于它的经典表达式
$m\boldsymbol{v}$，因而运动方程（17.5）变为
\begin{equation}
  m\frac{\mathrm{d}\boldsymbol{v}}{\mathrm{d}t}
  =e\boldsymbol{E}+\frac{e}{c}\boldsymbol{v}\times\boldsymbol{H}.
\end{equation}

下面我们还要导出粒子的动能随时间的变化率的方程\footnote{这里“动能”是指能量（9.4），它包含静能.}，即确定下面导数的方程，
\begin{equation*}
  \frac{\mathrm{d}\mathcal{E}_\mathrm{kin}}{\mathrm{d}t}
  =\frac{\mathrm{d}}{\mathrm{d}t}
  \left(\frac{mc^2}{\sqrt{1-\dfrac{v^2}{c^2}}}\right).
\end{equation*}

很容易验证
\begin{equation*}
  \frac{\mathrm{d}\mathcal{E}_\mathrm{kin}}{\mathrm{d}t}
  =\boldsymbol{v}\cdot\frac{\mathrm{d}\boldsymbol{p}}{\mathrm{d}t}.
\end{equation*}

将 (17.5) 式中的 $\mathrm{d}\boldsymbol{p}/\mathrm{d}t$ 代入，并注意到
$\boldsymbol{v}\times\boldsymbol{H}\cdot\boldsymbol{v}=0$，我们有
\begin{equation}
  \frac{\mathrm{d}\mathcal{E}_\mathrm{kin}}{\mathrm{d}t}=e\boldsymbol{E}\cdot\boldsymbol{v}.
\end{equation}

动能随时间的变化率就是场在单位时间内对粒子所做的功.从 (17.7) 可以看出，这个功等于速度与电场作用于电荷上的力的乘积.场在时间 $\mathrm{d}t$ 内所做的功，即在电荷移动
$\mathrm{d}\boldsymbol{r}$ 距离时所做的功，显然等于 $e\boldsymbol{E}\cdot\mathrm{d}\boldsymbol{r}$.

我们强调这个事实，即对电荷做功的仅仅是电场；磁场不能对在其中运动的电荷做功.这是因为磁场对电荷的作用力永与电荷的速度垂直.

力学方程相对于时间变号，即相对于将来与过去对调来说是不变的.换句话说，在力学中两个时间方向是等价的.这就意味着，假如某一种运动能按照力学方程进行，那么，相反的运动也是可能的，在这个相反的运动中，力学体系经历同样的状态，但是次序相反.

很容易看出，这对于相对论中的电磁场也成立.但是在这种情况下，除了将时间 $t$ 变为 $-t$ 外，我们还必须改变磁场的符号.事实上，很容易看出，假如我们作下列代换：
\begin{equation}
  t\to-t,\qquad \boldsymbol{E}\to\boldsymbol{E},\qquad
  \boldsymbol{H}\to-\boldsymbol{H},
\end{equation}
运动方程（17.5）是不变的.

按照 (17.3) 及 (17.4)，这并不改变标势，但是矢势变了号：
\begin{equation}
  \varphi\to\varphi,\qquad \boldsymbol{A}\to-\boldsymbol{A}.
\end{equation}

因此，假如某种运动在电磁场中是可能的，那么，相反的运动在 $\boldsymbol{H}$ 方向相反的电磁场中也是可能的.

\begin{xiti}

\noindent{\heiti 习\ 题}\quad 用粒子的速度及电场强度和磁场强度来表示它的加速度.

解：将 $\boldsymbol{p}=\boldsymbol{v}\mathcal{E}_\mathrm{kin}/c^2$ 代入（17.5），并按 (17.7) 表示
$\mathrm{d}\mathcal{E}_\mathrm{kin}/\mathrm{d}t$，求得
\begin{equation*}
  \dot{\boldsymbol{v}}=\frac{e}{m}\sqrt{1-\frac{v^2}{c^2}}
  \left\{\boldsymbol{E}+\frac{1}{c}\boldsymbol{v}\times\boldsymbol{H}
  -\frac{1}{c^2}\boldsymbol{v}\left(\boldsymbol{v}\cdot\boldsymbol{E}\right)\right\}.
\end{equation*}

\end{xiti}

\section{规范不变性}

现在让我们来研究场的势可唯一地确定到什么程度.首先，我们应该注意这个事实，就是场是由它对其内电荷运动所产生的影响来刻画的.但是在运动方程 (17.5) 内并未出现势，而只出现了场强 $\boldsymbol{E}$ 及 $\boldsymbol{H}$.所以两个场如果以相同两个矢量 $\boldsymbol{E}$ 及 $\boldsymbol{H}$ 来描述，那么，这两个场在物理上是全同的.

假如给定了势 $\boldsymbol{A}$ 及 $\varphi$，那么，根据 (17.3) 及 (17.4)，$\boldsymbol{E}$ 及
$\boldsymbol{H}$ 就由它们完全唯一地确定了.但是同一个场可以相应于不同的势.为了证明这件事，我们可将
$\partial f/\partial x^k$ 这个量加到势的每一个分量上，其中 $f$ 是坐标与时间的任意函数.因此，势 $A_k$ 变为
\begin{equation}
  A_k'=A_k-\frac{\partial f}{\partial x^k}.
\end{equation}

经过这样的改变，在作用量积分（16.1）中将出现附加项
\begin{equation}
  \frac{e}{c}\frac{\partial f}{\partial x^k}\mathrm{d}x^k
  =\mathrm{d}\left(\frac{e}{c}f\right),
\end{equation}
但是，将一个全微分加在作用量积分的被积函数内，运动方程并不受影响（见本教程第一卷§2）.

如果我们引入标势及矢势来代替四维势，而且用 $ct$，$x$，$y$，$z$ 来代替 $x^i$，那么（18.1）中的四个方程就可以写成下面的形式：
\begin{equation}
  \boldsymbol{A}'=\boldsymbol{A}+\operatorname{grad}f,\qquad
  \varphi'=\varphi-\frac{1}{c}\frac{\partial f}{\partial t}.
\end{equation}

很容易验证，当用 (18.3) 式定义的 $\boldsymbol{A}'$ 及 $\varphi'$ 代替 $\boldsymbol{A}$ 及
$\varphi$ 时，由方程 (17.3) 及 (17.4) 所定义的电场及磁场实际上并不改变.因此，势的变换 (18.1) 并不改变场.所以势并没有被唯一地确定，确定矢势仅仅精确到一个任意函数的梯度，而确定标势仅仅精确到同一个任意函数的时间导数.

特别是，我们显然可以加一个任意的常矢量到矢势上，加一个任意常数到标势上.这也可直接由如下事实看出来，$\boldsymbol{E}$ 及 $\boldsymbol{H}$ 仅仅包含 $\boldsymbol{A}$ 及 $\varphi$ 的导数，所以加些常量到 $\boldsymbol{A}$ 及 $\varphi$ 上，并不影响场的强度.

只有那些对于势变换 (18.3) 为不变的量才有物理意义；特别是，所有方程在这个变换下必须是不变的.这种不变性称为规范不变性（德文为 Eichinvarianz，英文为 Gauge invariance）\footnote{我们强调指出，这与（18.2）中假设 $e$ 的不变性有关.因此，电动力学方程的规范不变性（见下）同电荷守恒彼此密切相关.}.

因为势缺乏唯一性，我们就有可能去选择它们，使它们满足我们所选择的附加条件.我们强调指出，我们能够令它们满足一个条件，这是因为我们可以任意选择 (18.3) 中的函数 $f$.特别而言，我们总能这样来选择势，使得标势 $\varphi$ 为零.假如矢势不是零，那么，一般来说，我们不可能使它为零，因为条件 $\boldsymbol{A}=0$ 表示三个附加条件（即 $\boldsymbol{A}$ 的三个分量为零）.

\section{恒定电磁场}

所谓恒定电磁场就是与时间无关的电磁场.显然，恒定电磁场的势可以这样选择：使它们仅仅是坐标的函数，而不是时间的函数.恒定磁场同以前一样等于
$\boldsymbol{H}=\operatorname{rot}\boldsymbol{A}$.恒定电场等于
\begin{equation}
  \boldsymbol{E}=-\operatorname{grad}\varphi.
\end{equation}

因此，恒定电场仅为标势所决定，而恒定磁场仅为矢势所决定.

我们在前一节中已经知道，势并未被唯一地确定.但是很容易相信，假如我们用与时间无关的势来描述恒定电磁场，那么，我们可以在标势上加一个任意常数而不改变场（这个任意常数既与坐标无关，也与时间无关）.通常要给 $\varphi$ 加上一个附加条件，即在空间内某一特定点有一定的值；常常规定 $\varphi$ 在无穷远处的值为零.这时，上面所说的任意常数就被确定了，因而恒定场的标势也就被唯一地确定了.

另一方面，同以前一样，即使对于恒定电磁场，矢势也还是没有被唯一地确定；就是说，我们可以把坐标的一个任意函数的梯度加到矢势上.

现在我们要确定一个电荷在恒定电磁场内的能量.既然场是恒定的，那么，电荷的拉格朗日量也就不是时间的显函数.如我们所知道的，在这种情况下，能量是守恒的，而且它就是哈密顿量.

按照（16.6），我们得到
\begin{equation}
  \mathcal{E}=\frac{mc^2}{\sqrt{1-\dfrac{v^2}{c^2}}}+e\varphi.
\end{equation}

因之，由于场的存在，粒子的能量加上了一项 $e\varphi$，它是电荷在场内的势能.我们应该注意一个重要的事实，即能量仅与标势有关，而与矢势无关.这就意味着磁场不影响电荷的能量.只有电场才改变粒子的能量.这同如下事实有关，即磁场与电场不同，它对电荷不做功.

假如场强在空间所有点上都一样，那么这样的场称为均匀场.让我们用场强 $\boldsymbol{E}$ 来表示一个均匀电场的标势.很容易证明，对于一个均匀场来说，
\begin{equation}
  \varphi=-\boldsymbol{E}\cdot\boldsymbol{r}.
\end{equation}

实际上，因为 $\boldsymbol{E}$ 是常量，所以
$\nabla\left(\boldsymbol{E}\cdot\boldsymbol{r}\right)
=\left(\boldsymbol{E}\cdot\nabla\right)\boldsymbol{r}=\boldsymbol{E}$.

均匀磁场的矢势可以用场强 $\boldsymbol{H}$ 表示为：
\begin{equation}
  \boldsymbol{A}=\frac{1}{2}\boldsymbol{H}\times\boldsymbol{r}.
\end{equation}

事实上，我们注意到 $\boldsymbol{H}=\mathrm{const}$，利用矢量分析中熟知的公式可得
\begin{equation*}
  \operatorname{rot}\left(\boldsymbol{H}\times\boldsymbol{r}\right)
  =\boldsymbol{H}\operatorname{div}\boldsymbol{r}
  -\left(\boldsymbol{H}\cdot\nabla\right)\boldsymbol{r}
  =2\boldsymbol{H}
\end{equation*}
（注意，$\operatorname{div}\boldsymbol{r}=3$）.

均匀磁场的矢势也可以选择为下列形式：
\begin{equation}
  A_x=-Hy,\qquad A_y=A_z=0.
\end{equation}
（选择 $z$ 轴沿着 $\boldsymbol{H}$ 的方向）.很容易证明，这样选择了 $\boldsymbol{A}$，我们就有
$\boldsymbol{H}=\operatorname{rot}\boldsymbol{A}$.按照变换式 (18.3)，势 (19.4) 和 (19.5)
彼此之间的差为某个函数的梯度：由 (19.4) 加上 $\nabla f$ 就得到（19.5），式中 $f=-xyH/2$.

\begin{xiti}

\noindent{\heiti 习\ 题}\quad 在相对论力学中，写出恒定电磁场内粒子轨道的变分原理（即莫培督原理）.

解：莫培督原理可表述如下：假如一个粒子的能量是守恒的（在一个恒定场内运动），那么，它的轨道可从下面的变分方程确定：
\begin{equation*}
  \delta\int\boldsymbol{P}\cdot\mathrm{d}\boldsymbol{r}=0,
\end{equation*}
其中 $\boldsymbol{P}$ 是粒子的广义动量，这个广义动量 $\boldsymbol{P}$ 可以用能量及坐标的微分来表示；至于积分就应该沿粒子的轨道进行（见本教程第一卷§44）.将
$\boldsymbol{P}=\boldsymbol{p}+(e/c)\boldsymbol{A}$ 代入，注意 $\boldsymbol{p}$ 与
$\mathrm{d}\boldsymbol{r}$ 同方向，我们就得到
\begin{equation*}
  \delta\int\left(p\,\mathrm{d}l+\frac{e}{c}\boldsymbol{A}\cdot\mathrm{d}\boldsymbol{r}\right)=0,
\end{equation*}
其中 $\mathrm{d}l=\sqrt{\mathrm{d}\boldsymbol{r}^2}$ 是弧元.从
$p^2+m^2c^2=(\mathcal{E}-e\varphi)^2/c^2$ 来定 $p$，我们最后得到
\begin{equation*}
  \delta\int\left\{\sqrt{\left(\frac{\mathcal{E}-e\varphi}{c}\right)^2-m^2c^2}\,\mathrm{d}l
  +\frac{e}{c}\boldsymbol{A}\cdot\mathrm{d}\boldsymbol{r}\right\}=0.
\end{equation*}

\end{xiti}

\section{在恒定均匀电场中的运动}

现在我们研究电荷 $e$ 在均匀的恒定电场 $\boldsymbol{E}$ 中的运动.我们取电场的方向为 $x$ 轴.该运动显然在一个平面内进行.我们取这个平面为 $xy$ 平面.这时，运动方程（17.5）变为
\begin{equation*}
  \dot{p}_x=eE,\qquad \dot{p}_y=0
\end{equation*}
（上面的一点表示对时间 $t$ 微分），所以：
\begin{equation}
  p_x=eEt,\qquad p_y=p_0.
\end{equation}

我们把时间的参考点选在 $p_x=0$ 的时刻，$p_0$ 表示粒子在该时刻的动量.

粒子的动能（场内势能以外的能量）是
$\mathcal{E}_\mathrm{kin}=c\sqrt{m^2c^2+p^2}$，将（20.1）代入，我们求得，在现在这种情况下，
\begin{equation}
  \mathcal{E}_\mathrm{kin}
  =\sqrt{m^2c^4+c^2p_0^2+(ceEt)^2}
  =\sqrt{\mathcal{E}_0^2+(ceEt)^2},
\end{equation}
其中，$\mathcal{E}_0$ 是 $t=0$ 时的能量.

按照（9.8）式，粒子的速度 $\boldsymbol{v}=\boldsymbol{p}c^2/\mathcal{E}_\mathrm{kin}$.因而，对于速度
$v_x=\dot{x}$，我们有
\begin{equation*}
  \frac{\mathrm{d}x}{\mathrm{d}t}
  =\frac{p_xc^2}{\mathcal{E}_\mathrm{kin}}
  =\frac{c^2eEt}{\sqrt{\mathcal{E}_0^2+(ceEt)^2}}.
\end{equation*}

经过积分，我们求得
\begin{equation}
  x=\frac{1}{eE}\sqrt{\mathcal{E}_0^2+(ceEt)^2}
\end{equation}
（其中我们已经令积分常数等于零）\footnote{这个结果（对于 $p_0=0$）同具有恒定“固有加速度”$w_0=eE/m$ 的相对论运动问题的解一致（见 §7 的习题）.对于目前的情形，加速度的恒定性与速度 $\boldsymbol{V}$ 沿着场的方向时，洛伦兹变换不改变电场这一事实有关（见 §24）.}.

为了求得 $y$，我们有
\begin{equation*}
  \frac{\mathrm{d}y}{\mathrm{d}t}
  =\frac{p_yc^2}{\mathcal{E}_\mathrm{kin}}
  =\frac{p_0c^2}{\sqrt{\mathcal{E}_0^2+(ceEt)^2}},
\end{equation*}
从而
\begin{equation}
  y=\frac{p_0c}{eE}\operatorname{arsinh}\frac{ceEt}{\mathcal{E}_0}.
\end{equation}

利用 (20.4) 式通过 $y$ 来表示 $t$，并将它代入 (20.3)，我们得到轨道方程如下：
\begin{equation}
  x=\frac{\mathcal{E}_0}{eE}\cosh\frac{eEy}{p_0c}.
\end{equation}

由此可见，在均匀电场中一个电荷沿着悬链线运动.

假如粒子的速度 $v\ll c$，那么，我们可以使 $p_0=mv_0$，$\mathcal{E}_0=mc^2$，并将 (20.5)
式展开为 $1/c$ 的幂级数.略去高阶项，我们得到
\begin{equation*}
  x=\frac{eE}{2mv_0^2}y^2+\mathrm{const},
\end{equation*}
就是说，电荷沿着抛物线运动，这是我们在经典力学中熟知的结果.

\section{在恒定均匀磁场中的运动}

我们现在来研究电荷 $e$ 在恒定均匀磁场 $\boldsymbol{H}$ 中的运动.我们选择磁场的方向为 $z$ 轴的方向.我们将动量（(9.8) 式）
\begin{equation*}
  \boldsymbol{p}=\frac{\mathcal{E}\boldsymbol{v}}{c^2}
\end{equation*}
（式中 $\mathcal{E}$ 是粒子的能量，它在磁场中是恒定的）代入运动方程
\begin{equation*}
  \dot{\boldsymbol{p}}=\frac{e}{c}\boldsymbol{v}\times\boldsymbol{H},
\end{equation*}
就可将其改写成另一种形式.这样一来，运动方程就化成了下面的形式：
\begin{equation}
  \frac{\mathcal{E}}{c^2}\frac{\mathrm{d}\boldsymbol{v}}{\mathrm{d}t}
  =\frac{e}{c}\boldsymbol{v}\times\boldsymbol{H},
\end{equation}
或用分量表示为
\begin{equation}
  \dot{v}_x=\omega v_y,\qquad
  \dot{v}_y=-\omega v_x,\qquad
  \dot{v}_z=0,
\end{equation}
其中引入了符号
\begin{equation}
  \omega=\frac{ecH}{\mathcal{E}}.
\end{equation}

将（21.2）中的第二个方程乘以 $\mathrm{i}$，加到第一个方程中，就得到：
\begin{equation*}
  \frac{\mathrm{d}}{\mathrm{d}t}\left(v_x+\mathrm{i}v_y\right)
  =-\mathrm{i}\omega\left(v_x+\mathrm{i}v_y\right),
\end{equation*}
因而
\begin{equation*}
  v_x+\mathrm{i}v_y=a\,\mathrm{e}^{-\mathrm{i}\omega t},
\end{equation*}
其中 $a$ 是一个复常数.这个复常数可以写成 $a=v_{0t}\mathrm{e}^{-\mathrm{i}\alpha}$ 这样的形式，其中
$v_{0t}$ 及 $\alpha$ 都是实数.因此，
\begin{equation*}
  v_x+\mathrm{i}v_y=v_{0t}\,\mathrm{e}^{-\mathrm{i}(\omega t+\alpha)}.
\end{equation*}

将实部与虚部分开，我们得到
\begin{equation}
  v_x=v_{0t}\cos(\omega t+\alpha),\qquad
  v_y=-v_{0t}\sin(\omega t+\alpha).
\end{equation}

常数 $v_{0t}$ 及 $\alpha$ 都要由初始条件来决定，$\alpha$ 是初相位，至于 $v_{0t}$，则从（21.4）看出
\begin{equation*}
  v_{0t}=\sqrt{v_x^2+v_y^2},
\end{equation*}
即是说，$v_{0t}$ 是粒子在 $xy$ 平面内的速度，而且它在整个运动过程中保持不变.再次积分（21.4），我们得到
\begin{equation}
  x=x_0+r\sin(\omega t+\alpha),\qquad
  y=y_0+r\cos(\omega t+\alpha),
\end{equation}
其中
\begin{equation}
  r=\frac{v_{0t}}{\omega}=\frac{v_{0t}\mathcal{E}}{ecH}=\frac{cp_t}{eH}
\end{equation}
（$\boldsymbol{p}_t$ 是动量在 $xy$ 平面上的投影）.从 (21.2) 的第三个方程，我们得到
$v_z=v_{0z}$ 以及
\begin{equation}
  z=z_0+v_{0z}t.
\end{equation}

从（21.5）及（21.7），可以很明显地看出，电荷在均匀磁场中沿着螺旋线运动，螺旋线的轴沿着磁场的方向，螺旋线的半径由 (21.6) 式来决定.粒子的速度是常数.在
$v_{0z}=0$，即粒子没有沿磁场方向的速度分量的特殊情况下，粒子就在与磁场垂直的平面内作圆周运动.

从上面各公式我们可以看出，$\omega$ 是粒子在与磁场垂直的平面内转动的角频率.

假如粒子的速度很低，那么，我们就可以近似地令 $\mathcal{E}=mc^2$.这时频率 $\omega$ 变为
\begin{equation}
  \omega=\frac{eH}{mc}.
\end{equation}

现在我们假设磁场保持均匀，但数值和方向缓慢变化的情形，探究此时带电粒子的运动如何改变.

我们知道：在运动条件缓慢变化的情况下，某些被称为“浸渐不变量”的量保持恒定.因为在与磁场垂直的平面内的运动是周期性的，所以积分
\begin{equation*}
  I=\frac{1}{2\pi}\oint\boldsymbol{P}_t\cdot\mathrm{d}\boldsymbol{r}
\end{equation*}
是浸渐不变量，这个积分应该在运动的整个周期内来进行.在我们所讨论的情形下，就是沿着圆周来进行（$\boldsymbol{P}_t$
是广义动量在这个平面上的投影）\footnote{见本教程第一卷 §49.一般说来，沿特定坐标 $q$ 的一个周期进行的积分 $\oint p\,\mathrm{d}q$ 是浸渐不变量.在本情形下，垂直于 $\boldsymbol{H}$ 的平面内的两个坐标周期一致，我们已经写出的积分 $I$ 是两个相应浸渐不变量之和.然而，这些不变量每一个单独说来并没有特殊意义，因为它们依赖于场的矢势（非唯一的）选择.由此得出的浸渐不变量的非唯一性反映了如下事实，当我们认为磁场在全空间均匀时，原则上不能决定由 $\boldsymbol{H}$ 改变所产生的电场，因为它实际上依赖于无穷远处的特定条件.}.将
$\boldsymbol{P}_t=\boldsymbol{p}_t+(e/c)\boldsymbol{A}$ 代入，就得到
\begin{equation*}
  I=\frac{1}{2\pi}\oint\boldsymbol{P}_t\cdot\mathrm{d}\boldsymbol{r}
  =\frac{1}{2\pi}\oint\boldsymbol{p}_t\cdot\mathrm{d}\boldsymbol{r}
  +\frac{e}{2\pi c}\oint\boldsymbol{A}\cdot\mathrm{d}\boldsymbol{r}.
\end{equation*}

在第一项中我们应注意 $\boldsymbol{p}_t$ 的绝对值是常数，其方向沿着 $\mathrm{d}\boldsymbol{r}$，对第二项应用斯托克斯定理，并写出
$\operatorname{rot}\boldsymbol{A}=\boldsymbol{H}$，则得：
\begin{equation*}
  I=rp_t-\frac{e}{2c}Hr^2,
\end{equation*}
其中 $r$ 是轨道半径\footnote{对于给定方向的 $\boldsymbol{H}$ 考察电荷沿轨道的运动方向，我们发现，如果沿 $\boldsymbol{H}$ 看去，它是逆时针的.所以第二项是负号.}.将 $r$ 的表达式（21.6）代入，可得
\begin{equation}
  I=\frac{cp_t^2}{2eH}.
\end{equation}

由上式可以看出，对于 $H$ 的缓慢变化，切向动量 $p_t$ 的变化近似与 $\sqrt{H}$ 成正比.

这个结果也适用于另一种情况，即粒子在并不严格均匀的磁场中沿螺旋轨道运动（使得场在可以同螺旋线的半径和步长相比的距离上变化很小）.这样的运动可以看成在随时间移动的圆轨道上进行，场相对于轨道表现为随时间变化但却保持均匀.于是我们可以说，垂直于场方向的动量分量按照规律
$p_t=\sqrt{CH}$ 变化，这里 $C$ 是一常数，$H$ 是坐标的给定函数.另一方面，正如任何恒定磁场中的运动那样，粒子的能量（因而其动量的平方
$p^2$）保持不变.因此动量的纵向分量按照如下公式变化：
\begin{equation}
  p_l^2=p^2-p_t^2=p^2-CH(x,y,z).
\end{equation}

因为我们总是应当有 $p_l^2\geqslant0$，粒子穿入场足够强的区域（$CH>p^2$）是不可能的.在沿场增加方向的运动进程中，螺旋轨道的半径按
$p_t/H$ 成比例地减小（即比例于 $1/\sqrt{H}$），而步长比例于 $p_l$.在达到 $p_l$ 变为零的边界处，粒子被反射：在继续沿相同方向转动的同时，开始朝场梯度相反的方向运动.

场的不均匀性也导致另一种现象——粒子螺旋轨道的导引中心（像称圆轨道的中心那样）缓慢地横向移动（漂移）；下一节习题3将讨论这个问题.

\begin{xiti}

\noindent{\heiti 习\ 题}\quad 将一个带电的空间振子置于均匀恒定磁场内；这个振子（在没有置于场内时的振动）的本征频率是
$\omega_0$，求它在置于场内后的振动频率.

解：振子在磁场（磁场方向沿着 $z$ 轴）内的受迫振动方程是
\begin{equation*}
  \ddot{x}+\omega_0^2x=\frac{eH}{mc}\dot{y},\qquad
  \ddot{y}+\omega_0^2y=-\frac{eH}{mc}\dot{x},\qquad
  \ddot{z}+\omega_0^2z=0.
\end{equation*}

用 $\mathrm{i}$ 乘第二个方程并与第一个方程相加，得到
\begin{equation*}
  \ddot{\xi}+\omega_0^2\xi=-\mathrm{i}\frac{eH}{mc}\dot{\xi},
\end{equation*}
其中 $\xi=x+\mathrm{i}y$.由此可以得到振子在与磁场垂直的平面内的振动频率为
\begin{equation*}
  \omega=\sqrt{\omega_0^2+\frac{1}{4}\left(\frac{eH}{mc}\right)^2}
  \pm\frac{eH}{2mc}.
\end{equation*}

假如磁场 $H$ 很弱，这个公式就变为
\begin{equation*}
  \omega=\omega_0\pm\frac{eH}{2mc}.
\end{equation*}

沿着场的方向的振动将保持不变.

\end{xiti}

\section{电荷在均匀恒定的电场和磁场中的运动}

最后，我们来研究在电场及磁场都存在并且都是均匀恒定的情况下，一个电荷如何运动.我们的讨论只限于粒子的速度
$v\ll c$ 的情形，因此质点的动量 $\boldsymbol{p}=m\boldsymbol{v}$；以后我们将要知道，出现这种情形的必要条件是电场比磁场小得多.

我们选择 $\boldsymbol{H}$ 的方向为 $z$ 轴的方向，而选择通过 $\boldsymbol{H}$ 及 $\boldsymbol{E}$
的平面为 $yz$ 平面.这时，运动方程
\begin{equation*}
  m\dot{\boldsymbol{v}}=e\boldsymbol{E}
  +\frac{e}{c}\boldsymbol{v}\times\boldsymbol{H}
\end{equation*}
可以写成如下式：
\begin{equation}
\begin{aligned}
  &m\ddot{x}=\frac{e}{c}\dot{y}H,\\
  &m\ddot{y}=eE_y-\frac{e}{c}\dot{x}H,\\
  &m\ddot{z}=eE_z.
\end{aligned}
\end{equation}

由上面的第三个方程，我们可以看出，电荷以匀加速度沿着 $z$ 轴方向运动，就是说，
\begin{equation}
  z=\frac{eE_z}{2m}t^2+v_{0z}t.
\end{equation}

用 $\mathrm{i}$ 乘（22.1）中的第二个方程，再与第一个方程联立，我们得到
\begin{equation*}
  \frac{\mathrm{d}}{\mathrm{d}t}\left(\dot{x}+\mathrm{i}\dot{y}\right)
  +\mathrm{i}\omega\left(\dot{x}+\mathrm{i}\dot{y}\right)
  =\mathrm{i}\frac{e}{m}E_y
\end{equation*}
（$\omega=eH/mc$）.将 $\dot{x}+\mathrm{i}\dot{y}$ 当做未知量，上面方程的解就等于上面的方程略去等号右边项的解，与该方程保留等号右边项的一个特解之和.第一个解是
$a\,\mathrm{e}^{-\mathrm{i}\omega t}$，第二个解是 $eE_y/m\omega=cE_y/H$，因此，
\begin{equation*}
  \dot{x}+\mathrm{i}\dot{y}=a\,\mathrm{e}^{-\mathrm{i}\omega t}+\frac{cE_y}{H}.
\end{equation*}

常数 $a$ 一般来说是个复数.将它写成 $a=b\,\mathrm{e}^{\mathrm{i}\alpha}$ 的形式，其中 $b$ 及
$\alpha$ 为实数，我们可以看出，既然 $a$ 被 $\mathrm{e}^{-\mathrm{i}\omega t}$ 乘了，那么，只要我们选择时间原点得当，就可以赋予相位
$\alpha$ 以任何一个值.我们适当选择时间原点，使 $\alpha$ 为实数.将
$\dot{x}+\mathrm{i}\dot{y}$ 分解为实部及虚部，我们便得到
\begin{equation}
  \dot{x}=a\cos\omega t+c\frac{E_y}{H},\qquad
  \dot{y}=-a\sin\omega t.
\end{equation}

在 $t=0$ 时，速度沿着 $x$ 轴.

我们可以看出，粒子的速度是时间的周期函数.它们的平均值是：
\begin{equation*}
  \overline{\dot{x}}=\frac{cE_y}{H},\qquad \overline{\dot{y}}=0.
\end{equation*}

电荷在正交的电场和磁场中运动的这个平均速度常称为电漂移速度.它的方向与两个场都垂直并与电荷的正负无关.可以把它写成矢量形式：
\begin{equation}
  \overline{\boldsymbol{v}}=\frac{c\,\boldsymbol{E}\times\boldsymbol{H}}{H^2}.
\end{equation}

这一节的所有公式，都假设了粒子的速度比光速小得多；可以看出，为了实现这种情形，特别要求电场与磁场必须满足下面的条件：
\begin{equation}
  \frac{E_y}{H}\ll1,
\end{equation}
而 $E_y$ 及 $H$ 的绝对大小可以是任意的.

将方程（22.3）再积分一次，并这样来选择积分常数，使当 $t=0$ 时，$x=y=0$，我们就可得到
\begin{equation}
\begin{aligned}
  &x=\frac{a}{\omega}\sin\omega t+\frac{cE_y}{H}t,\\
  &y=\frac{a}{\omega}\left(\cos\omega t-1\right).
\end{aligned}
\end{equation}

将以上二式看做一个曲线的参数方程，这两个方程定义一条次摆线.至于轨道在 $xy$
平面上的投影到底是如图6a还是如图6b所示的，那就得看 $a$ 的绝对值是大于还是小于 $cE_y/H$.

\par\nopagebreak\noindent\begin{minipage}{\linewidth}\centering\includegraphics[width=5cm]{fig6.pdf}\\[0.3em]{\small 图 6}\end{minipage}\par\nopagebreak

假如 $a=-cE_y/H$，那么，(22.6) 就变为
\begin{equation}
\begin{aligned}
  &x=\frac{cE_y}{\omega H}\left(\omega t-\sin\omega t\right),\\
  &y=\frac{cE_y}{\omega H}\left(1-\cos\omega t\right),
\end{aligned}
\end{equation}
就是说，轨道在 $xy$ 平面上的投影是一条旋轮线（图6c）.

\begin{xiti}

\noindent{\heiti 习题1}\quad 求电荷在平行均匀电场和磁场中的相对论运动.

解：磁场对沿 $\boldsymbol{E}$ 和 $\boldsymbol{H}$ 共同方向（$z$ 轴）上的运动没有影响，因而 $z$
方向上的运动只在电场的影响下发生；于是根据§20，我们得到：
\begin{equation*}
  z=\frac{\mathcal{E}_\mathrm{kin}}{eE},\qquad
  \mathcal{E}_\mathrm{kin}=\sqrt{\mathcal{E}_0^2+(ceEt)^2}.
\end{equation*}

对于 $xy$ 平面上的运动，我们有方程
\begin{equation*}
  \dot{p}_x=\frac{e}{c}Hv_y,\qquad
  \dot{p}_y=-\frac{e}{c}Hv_x,
\end{equation*}
或
\begin{equation*}
  \frac{\mathrm{d}}{\mathrm{d}t}\left(p_x+\mathrm{i}p_y\right)
  =-\mathrm{i}\frac{eH}{c}\left(v_x+\mathrm{i}v_y\right)
  =-\frac{\mathrm{i}eHc}{\mathcal{E}_\mathrm{kin}}\left(p_x+\mathrm{i}p_y\right).
\end{equation*}

由此得
\begin{equation*}
  p_x+\mathrm{i}p_y=p_t\,\mathrm{e}^{-\mathrm{i}\varphi},
\end{equation*}
式中 $\boldsymbol{p}_t$ 是动量在 $xy$ 平面上投影的恒定值，而辅助量 $\varphi$ 由如下关系决定
\begin{equation*}
  \mathrm{d}\varphi=eHc\,\frac{\mathrm{d}t}{\mathcal{E}_\mathrm{kin}},
\end{equation*}
由此
\begin{equation*}
  ct=\frac{\mathcal{E}_0}{eE}\sinh\frac{E}{H}\varphi.\tag{1}
\end{equation*}

然后我们有：
\begin{equation*}
  p_x+\mathrm{i}p_y=p_t\,\mathrm{e}^{-\mathrm{i}\varphi}
  =\frac{\mathcal{E}_\mathrm{kin}}{c^2}\left(\dot{x}+\mathrm{i}\dot{y}\right)
  =\frac{eH}{c}\frac{\mathrm{d}(x+\mathrm{i}y)}{\mathrm{d}\varphi},
\end{equation*}
所以
\begin{equation*}
  x=\frac{cp_t}{eH}\sin\varphi,\qquad
  y=\frac{cp_t}{eH}\cos\varphi.\tag{2}
\end{equation*}

公式 (1)，(2) 和公式
\begin{equation*}
  z=\frac{\mathcal{E}_0}{eE}\cosh\frac{E}{H}\varphi,\tag{3}
\end{equation*}
一起以参数形式决定了粒子的运动.轨道是一条螺旋线，半径为 $cp_t/eH$，步长单调增加，粒子沿着该螺旋线以减小的角速度
$\varphi=eHc/\mathcal{E}_\mathrm{kin}$ 运动，沿 $z$ 轴的速度趋向值 $c$.

\noindent{\heiti 习题2}\quad 求电荷在相互垂直且数值相等的电场和磁场中的相对论运动\footnote{相互垂直而数值不等的 $\boldsymbol{E}$ 和 $\boldsymbol{H}$ 场中的运动问题，可以通过适当的参考系变换，化为纯电场或纯磁场中运动的问题，见 §25.}.

解：选 $z$ 轴沿 $\boldsymbol{H}$ 方向，$y$ 轴沿 $\boldsymbol{E}$ 方向，令 $E=H$，我们写出运动方程：
\begin{equation*}
  \frac{\mathrm{d}p_x}{\mathrm{d}t}=\frac{e}{c}Ev_y,\qquad
  \frac{\mathrm{d}p_y}{\mathrm{d}t}=eE\left(1-\frac{v_x}{c}\right),\qquad
  \frac{\mathrm{d}p_z}{\mathrm{d}t}=0
\end{equation*}
以及——作为它们的结果——(17.7) 式：
\begin{equation*}
  \frac{\mathrm{d}\mathcal{E}_\mathrm{kin}}{\mathrm{d}t}=eEv_y.
\end{equation*}

由这些方程，我们有：
\begin{equation*}
  p_z=\mathrm{const},\qquad
  \mathcal{E}_\mathrm{kin}-cp_x=\mathrm{const}\equiv\alpha.
\end{equation*}

再用方程
\begin{equation*}
  \mathcal{E}_\mathrm{kin}^2-c^2p_x^2
  =\left(\mathcal{E}_\mathrm{kin}+cp_x\right)
  \left(\mathcal{E}_\mathrm{kin}-cp_x\right)
  =c^2p_y^2+\varepsilon^2
\end{equation*}
（式中 $\varepsilon^2=m^2c^4+c^2p_z^2=\mathrm{const}$）我们有：
\begin{equation*}
  \mathcal{E}_\mathrm{kin}+cp_x=\frac{1}{\alpha}\left(c^2p_y^2+\varepsilon^2\right),
\end{equation*}
所以
\begin{equation*}
  \mathcal{E}_\mathrm{kin}=\frac{\alpha}{2}
  +\frac{c^2p_y^2+\varepsilon^2}{2\alpha},
\end{equation*}
\begin{equation*}
  p_x=-\frac{\alpha}{2c}+\frac{c^2p_y^2+\varepsilon^2}{2\alpha c}.
\end{equation*}

我们进一步写出
\begin{equation*}
  \mathcal{E}_\mathrm{kin}\frac{\mathrm{d}p_y}{\mathrm{d}t}
  =eE\left(\mathcal{E}_\mathrm{kin}-\frac{\mathcal{E}_\mathrm{kin}v_x}{c}\right)
  =eE\left(\mathcal{E}_\mathrm{kin}-cp_x\right)=eE\alpha,
\end{equation*}
由此
\begin{equation*}
  2eEt=\left(1+\frac{\varepsilon^2}{\alpha^2}\right)p_y
  +\frac{c^2}{3\alpha^2}p_y^3.\tag{1}
\end{equation*}

为了决定轨道，我们在方程
\begin{equation*}
  \frac{\mathrm{d}x}{\mathrm{d}t}=\frac{c^2p_x}{\mathcal{E}_\mathrm{kin}},\ \cdots
\end{equation*}
中作变量代换，对变量 $p_y$ 使用关系式
$\mathrm{d}t=\mathcal{E}_\mathrm{kin}\,\mathrm{d}p_y/eE\alpha$，之后积分给出公式：
\begin{equation*}\tag{2}
\begin{aligned}
  &x=\frac{c}{2eE}\left(-1+\frac{\varepsilon^2}{\alpha^2}\right)p_y
  +\frac{c^3}{6\alpha^2eE}p_y^3,\\
  &y=\frac{c^2}{2\alpha eE}p_y^2,\qquad
  z=\frac{p_zc^2}{eE\alpha}p_y.
\end{aligned}
\end{equation*}

公式 (1) 和 (2) 以参数形式（参数 $p_y$）完全决定了粒子的运动.请注意如下事实，速度在垂直于 $\boldsymbol{E}$ 和
$\boldsymbol{H}$ 轴的方向（$x$ 轴）增加得最快.

\noindent{\heiti 习题3}\quad 试求非相对论带电粒子在准均匀磁场中轨道导引中心的漂移速度（H. Alfv\'en, 1940）.

解：我们首先假设粒子在圆轨道上运动，即它的速度没有（沿着场的）纵向分量.写出形如
$\boldsymbol{r}=\boldsymbol{R}(t)+\boldsymbol{\zeta}(t)$ 的轨道方程，式中 $\boldsymbol{R}(t)$
是导引中心的径矢（为时间的缓变函数），$\boldsymbol{\zeta}(t)$ 是一快速振荡的量，描述围绕导引中心的转动.我们在振荡的（圆）运动周期内对作用于粒子上的力
$(e/c)\dot{\boldsymbol{r}}\times\boldsymbol{H}(\boldsymbol{r})$ 进行平均（比较本教程第一卷§30）.将这个式子中的函数
$\boldsymbol{H}(\boldsymbol{r})$ 以 $\boldsymbol{\zeta}$ 的幂展开：
\begin{equation*}
  \boldsymbol{H}(\boldsymbol{r})=\boldsymbol{H}(\boldsymbol{R})
  +\left(\boldsymbol{\zeta}\cdot\nabla\right)\boldsymbol{H}(\boldsymbol{R}).
\end{equation*}

平均后，振荡量 $\boldsymbol{\zeta}(t)$ 的一次项为零，二次项产生一附加的力
\begin{equation*}
  \boldsymbol{f}=\frac{e}{c}\,
  \overline{\boldsymbol{\zeta}\times
  \left(\boldsymbol{\zeta}\cdot\nabla\right)\boldsymbol{H}}.
\end{equation*}

对于圆轨道
\begin{equation*}
  \dot{\boldsymbol{\zeta}}=\omega\,\boldsymbol{\zeta}\times\boldsymbol{n},\qquad
  \zeta=\frac{v_\perp}{\omega},
\end{equation*}
式中 $\boldsymbol{n}$ 是沿 $\boldsymbol{H}$ 的单位矢量；频率 $\omega=eH/mc$，$v_\perp$
是粒子圆周运动的速度.矢量 $\boldsymbol{\zeta}$ 在一平面内转动（该平面垂直于 $\boldsymbol{n}$），其分量乘积的平均值为：
\begin{equation*}
  \overline{\zeta_\alpha\zeta_\beta}=\frac{1}{2}\zeta^2\delta_{\alpha\beta},
\end{equation*}
式中 $\delta_{\alpha\beta}$ 是这个平面内的单位张量.结果我们得到：
\begin{equation*}
  \boldsymbol{f}=-\frac{mv_\perp^2}{2H}
  \left(\boldsymbol{n}\times\nabla\right)\times\boldsymbol{H}.
\end{equation*}

由于恒定场 $\boldsymbol{H}(\boldsymbol{R})$ 满足方程 $\operatorname{div}\boldsymbol{H}=0$ 和
$\operatorname{rot}\boldsymbol{H}=0$，我们有：
\begin{equation*}
\begin{aligned}
  &\left(\boldsymbol{n}\times\nabla\right)\times\boldsymbol{H}
  =-\boldsymbol{n}\operatorname{div}\boldsymbol{H}
  +\left(\boldsymbol{n}\cdot\nabla\right)\boldsymbol{H}
  +\boldsymbol{n}\times\operatorname{rot}\boldsymbol{H}
  =\left(\boldsymbol{n}\cdot\nabla\right)\boldsymbol{H}\\
  &=\boldsymbol{H}\left(\boldsymbol{n}\cdot\nabla\right)\boldsymbol{n}
  +\boldsymbol{n}\left(\boldsymbol{n}\cdot\nabla\boldsymbol{H}\right).
\end{aligned}
\end{equation*}

我们对垂直于 $\boldsymbol{n}$、产生轨道漂移的力感兴趣；它等于：
\begin{equation*}
  \boldsymbol{f}=-\frac{mv_\perp^2}{2}
  \left(\boldsymbol{n}\cdot\nabla\right)\boldsymbol{n}
  =\frac{mv_\perp^2}{2\rho}\,\boldsymbol{\nu},
\end{equation*}
式中 $\rho$ 是场的力线在给定点的曲率半径，$\boldsymbol{\nu}$ 是从曲率中心指向该点的单位矢量.

粒子也具有纵向速度 $v_\parallel$（沿着 $\boldsymbol{n}$）的情况可以化为上述情形，只需我们换到绕力线（导引中心的轨道）的瞬时曲率中心转动的参考系即可，转动的角速度为
$v_\parallel/\rho$.在这个参考系中，粒子没有纵向速度，但有一附加横向力，即离心力
$\boldsymbol{\nu}\,mv_\parallel^2/\rho$.因此，总的横向力是
\begin{equation*}
  \boldsymbol{f}_\perp=\boldsymbol{\nu}\,\frac{m}{\rho}
  \left(v_\parallel^2+\frac{v_\perp^2}{2}\right).
\end{equation*}

这个力等价于强度为 $\boldsymbol{f}_\perp/e$ 的恒定电场.按照 (22.4)，它引起轨道导引中心的漂移，速度为
\begin{equation*}
  v_{\mathrm{d}}=\frac{1}{\omega\rho}
  \left(v_\parallel^2+\frac{v_\perp^2}{2}\right)\boldsymbol{\nu}\times\boldsymbol{n}.
\end{equation*}

这个速度的正负号依赖于电荷的正负号.

\end{xiti}

\section{电磁场张量}

在§17中，我们从写成三维形式的拉格朗日量（16.4）导出了电荷在场内运动的方程.现在我们直接从写成四维形式的作用量（16.1）导出同样的方程.

最小作用量原理是：
\begin{equation}
  \delta S=\delta\int_a^b\left(-mc\,\mathrm{d}s-\frac{e}{c}A_i\,\mathrm{d}x^i\right)=0.
\end{equation}

注意到 $\mathrm{d}s=\sqrt{\mathrm{d}x_i\,\mathrm{d}x^i}$，我们便求得（为了简便起见，下面我们略去积分限 $a$ 和
$b$）：
\begin{equation*}
  \delta S=-\int\left(mc\,\frac{\mathrm{d}x_i\,\mathrm{d}\delta x^i}{\mathrm{d}s}
  +\frac{e}{c}A_i\,\mathrm{d}\delta x^i
  +\frac{e}{c}\delta A_i\,\mathrm{d}x^i\right)=0.
\end{equation*}

将被积函数中前两项作分部积分，并且在第一项中令 $\mathrm{d}x_i/\mathrm{d}s=u_i$，其中 $u_i$
是四维速度的分量.于是
\begin{equation}
  \int\left(mc\,\mathrm{d}u_i\,\delta x^i+\frac{e}{c}\delta x^i\,\mathrm{d}A_i
  -\frac{e}{c}\delta A_i\,\mathrm{d}x^i\right)
  -\left(mc\,u_i+\frac{e}{c}A_i\right)\delta x^i=0.
\end{equation}

由于积分的变分是在两个边界具有固定坐标值的条件下取的，上式中的第二项等于零，此外：
\begin{equation*}
  \delta A_i=\frac{\partial A_i}{\partial x^k}\delta x^k,\qquad
  \mathrm{d}A_i=\frac{\partial A_i}{\partial x^k}\mathrm{d}x^k,
\end{equation*}
因此
\begin{equation*}
  \int\left(mc\,\mathrm{d}u_i\,\delta x^i
  +\frac{e}{c}\frac{\partial A_i}{\partial x^k}\delta x^i\,\mathrm{d}x^k
  -\frac{e}{c}\frac{\partial A_i}{\partial x^k}\mathrm{d}x^i\,\delta x^k\right)=0.
\end{equation*}

在第一项中，我们写 $\mathrm{d}u_i=\dfrac{\mathrm{d}u_i}{\mathrm{d}s}\mathrm{d}s$；在第二和第三项中，我们写
$\mathrm{d}x^i=u^i\mathrm{d}s$；在第三项中将指标 $i$ 与 $k$ 交换（因为 $i$ 和 $k$
都是求和指标，所以交换以后什么也不改变）.于是，
\begin{equation*}
  \int\left[mc\,\frac{\mathrm{d}u_i}{\mathrm{d}s}
  -\frac{e}{c}\left(\frac{\partial A_k}{\partial x^i}
  -\frac{\partial A_i}{\partial x^k}\right)u^k\right]\delta x^i\,\mathrm{d}s=0.
\end{equation*}

由于 $\delta x^i$ 的任意性，我们可以推断，被积函数必须为零，即
\begin{equation*}
  mc\,\frac{\mathrm{d}u_i}{\mathrm{d}s}
  -\frac{e}{c}\left(\frac{\partial A_k}{\partial x^i}
  -\frac{\partial A_i}{\partial x^k}\right)u^k=0.
\end{equation*}

现在我们引入下面的符号：
\begin{equation}
  F_{ik}=\frac{\partial A_k}{\partial x^i}-\frac{\partial A_i}{\partial x^k}.
\end{equation}

反对称张量 $F_{ik}$ 称为电磁场张量.这样一来，运动方程取下面的形式：
\begin{equation}
  mc\,\frac{\mathrm{d}u^i}{\mathrm{d}s}=\frac{e}{c}F^{ik}u_k.
\end{equation}

这就是四维形式的电荷运动方程.

将 $A_i=(\varphi,-\boldsymbol{A})$ 的值代入定义式 (23.3)，我们很容易看出张量 $F_{ik}$
的各个分量的意义.结果可以写成一个矩阵，其中指标 $i=0,1,2,3$ 表示行，指标 $k$ 表示列：
\begin{equation}
  F_{ik}=
  \begin{pmatrix}
    0 & E_x & E_y & E_z\\
    -E_x & 0 & -H_z & H_y\\
    -E_y & H_z & 0 & -H_x\\
    -E_z & -H_y & H_x & 0
  \end{pmatrix},
  \qquad
  F^{ik}=
  \begin{pmatrix}
    0 & -E_x & -E_y & -E_z\\
    E_x & 0 & -H_z & H_y\\
    E_y & H_z & 0 & -H_x\\
    E_z & -H_y & H_x & 0
  \end{pmatrix}.
\end{equation}

上式可以更简洁地写成（见§6）：
\begin{equation*}
  F_{ik}=(\boldsymbol{E},\boldsymbol{H}),\qquad
  F^{ik}=(-\boldsymbol{E},\boldsymbol{H}).
\end{equation*}

由此可见，电场强度与磁场强度的分量是同一四维电磁场张量的分量.

改变到三维符号，容易验证，(23.4) 的三个空间分量 $(i=1,2,3)$ 与矢量运动方程（17.5）相同，而时间分量
$(i=0)$ 给出功方程 (17.7).后者是运动方程的推论；四个方程中只有三个独立这一事实，也容易通过直接将 (23.4)
两边乘以 $u^i$ 发现.由于四维矢量 $u^i$ 和 $\mathrm{d}u_i/\mathrm{d}s$ 的正交性，此时方程的左边为零，而由于
$F_{ik}$ 的反对称性，方程的右边为零.

假如在变分 $\delta S$ 中，只考虑真实的轨道，那么，(23.2) 中的第一项就恒等于零.这时，第二项（其中积分上限认为是变化的）给出作为坐标函数的作用量的微分.因之
\begin{equation}
  \delta S=-\left(mc\,u_i+\frac{e}{c}A_i\right)\delta x^i.
\end{equation}

由此可得，
\begin{equation}
  -\frac{\partial S}{\partial x^i}=mc\,u_i+\frac{e}{c}A_i
  =p_i+\frac{e}{c}A_i.
\end{equation}

四维矢量 $-\partial S/\partial x^i$ 是粒子的广义动量四维矢量 $P_i$.代入分量值 $p_i$ 和 $A_i$
我们求出：
\begin{equation}
  P_i=\left(\frac{\mathcal{E}_\mathrm{kin}+e\varphi}{c},\,
  \boldsymbol{p}+\frac{e}{c}\boldsymbol{A}\right),
\end{equation}
正如预期，这个四维矢量 $P_i$ 的三个空间分量构成三维广义动量矢量 (16.5)，其时间分量是
$\mathcal{E}/c$，其中 $\mathcal{E}$ 是电荷在场中的总能量.

\section{场的洛伦兹变换}

在本节内，我们将寻找场的变换公式，利用这些公式，假如在某一个惯性参考系内场是已知的，在另一个惯性参考系内，我们也能够决定这个场.

势的变换公式可以直接从四维矢量的变换通式 (6.1) 得出.只需记起
$A^i=(\varphi,\boldsymbol{A})$，我们就很容易得到
\begin{equation}
  \varphi=\frac{\varphi'+\dfrac{V}{c}A_x'}
  {\sqrt{1-\dfrac{V^2}{c^2}}},\qquad
  A_x=\frac{A_x'+\dfrac{V}{c}\varphi'}
  {\sqrt{1-\dfrac{V^2}{c^2}}},\qquad
  A_y=A_y',\qquad
  A_z=A_z'.
\end{equation}

反对称二阶张量（如 $F^{ik}$）的变换公式可以从§6的习题2找到：分量 $F^{23}$ 和 $F^{01}$
不变，而分量 $F^{02},F^{03}$ 和 $F^{12},F^{13}$ 分别像 $x^0$ 和 $x^1$
一样变换.按照 (23.5) 把 $F^{ik}$ 的分量用场 $\boldsymbol{E}$ 和 $\boldsymbol{H}$
的分量来表示，我们就得到电场的变换公式：
\begin{equation}
  E_x=E_x',\qquad
  E_y=\frac{E_y'+\dfrac{V}{c}H_z'}{\sqrt{1-\dfrac{V^2}{c^2}}},\qquad
  E_z=\frac{E_z'-\dfrac{V}{c}H_y'}{\sqrt{1-\dfrac{V^2}{c^2}}}
\end{equation}
和磁场的变换公式：
\begin{equation}
  H_x=H_x',\qquad
  H_y=\frac{H_y'-\dfrac{V}{c}E_z'}{\sqrt{1-\dfrac{V^2}{c^2}}},\qquad
  H_z=\frac{H_z'+\dfrac{V}{c}E_y'}{\sqrt{1-\dfrac{V^2}{c^2}}}.
\end{equation}

因此，电场与磁场，正如大多数的物理量一样，是相对的；就是说，它们在不同的参考系中有不同的特性.特别是，电场或磁场在一个参考系中可以等于零，而同时在另外一个参考系中却又存在.

公式（24.2）和（24.3）在 $V\ll c$
的情况下将大大简化.精确到 $V/c$ 的数量级，我们有
\begin{equation*}
\begin{aligned}
  &E_x=E_x',\qquad
  E_y=E_y'+\frac{V}{c}H_z',\qquad
  E_z=E_z'-\frac{V}{c}H_y';\\
  &H_x=H_x',\qquad
  H_y=H_y'-\frac{V}{c}E_z',\qquad
  H_z=H_z'+\frac{V}{c}E_y'.
\end{aligned}
\end{equation*}

这些公式可以写成矢量形式
\begin{equation}
  \boldsymbol{E}=\boldsymbol{E}'
  +\frac{1}{c}\boldsymbol{H}'\times\boldsymbol{V},\qquad
  \boldsymbol{H}=\boldsymbol{H}'
  -\frac{1}{c}\boldsymbol{E}'\times\boldsymbol{V}.
\end{equation}

从 $K'$ 到 $K$ 的逆变换公式可以通过改变 $V$ 的符号并移动撇号直接由
(24.2)—(24.4) 求得.

假如在 $K'$ 系中磁场 $\boldsymbol{H}'=0$，那么，根据（24.2）及（24.3），我们可以很容易验证，在 $K$
系中，电场与磁场之间存在着下面的关系：
\begin{equation}
  \boldsymbol{H}=\frac{1}{c}\boldsymbol{V}\times\boldsymbol{E}.
\end{equation}

假如在 $K'$ 系中，$\boldsymbol{E}'=0$，那么，在 $K$ 系中
\begin{equation}
  \boldsymbol{E}=-\frac{1}{c}\boldsymbol{V}\times\boldsymbol{H}.
\end{equation}

因此，在以上两种情况下，电场与磁场在 $K$ 系中都是相互垂直的.

反向应用这些公式也是有意义的：如果场 $\boldsymbol{E}$ 和 $\boldsymbol{H}$
在某个参考系 $K$ 中相互垂直（但数值不等），那就存在一个参考系 $K'$，其中场是纯电场或纯磁场.

这个系统的速度 $\boldsymbol{V}$（相对于 $K$）垂直于 $\boldsymbol{E}$ 和 $\boldsymbol{H}$，并在数值上等于
$cH/E$（这种情况必须有 $H<E$），或者等于 $cE/H$（$E<H$ 的情形）.

\section{场的不变量}

从电磁场张量我们可以造出一些不变量，这些不变量在从一个惯性参考系过渡到另一个惯性参考系时保持不变.

从场的四维表示出发，用反对称四维张量 $F^{ik}$，容易得出这些不变量的形式.显然我们可以从这个张量的分量构造如下不变量：
\begin{equation}
  F_{ik}F^{ik}=\mathrm{inv},
\end{equation}
\begin{equation}
  e^{iklm}F_{ik}F_{lm}=\mathrm{inv},
\end{equation}
式中 $e^{iklm}$ 是四阶全反对称单位张量（见§6）.第一个量是标量，而第二个量是赝标量（张量 $F^{ik}$
同其对偶张量的乘积）\footnote{我们注意到，赝标量（25.2）也可以表示为一个四维散度：
\begin{equation*}
  e^{iklm}F_{ik}F_{lm}=4\frac{\partial}{\partial x^i}
  \left(e^{iklm}A_k\frac{\partial}{\partial x^l}A_m\right),
\end{equation*}
这一点可以通过 $e^{iklm}$ 的反对称性容易得到验证.}.

用 (23.5) 借助 $\boldsymbol{E}$ 和 $\boldsymbol{H}$ 的分量来表示 $F^{ik}$，容易证明，在三维形式下，这些不变量为：
\begin{equation}
  H^2-E^2=\mathrm{inv},
\end{equation}
\begin{equation}
  \boldsymbol{E}\cdot\boldsymbol{H}=\mathrm{inv}.
\end{equation}

其中第二个的赝标量特征从如下事实可以看出，它是极矢量 $\boldsymbol{E}$ 和轴矢量 $\boldsymbol{H}$
的乘积（尽管其平方 $(\boldsymbol{E}\cdot\boldsymbol{H})^2$ 是一个真标量）.

从上面得到的两个表达式的不变性，我们得出如下定理.如果电场和磁场在任一惯性系中相互垂直，即
$\boldsymbol{E}\cdot\boldsymbol{H}=0$，那么它们在其他每个惯性系中也垂直.如果 $\boldsymbol{E}$ 和
$\boldsymbol{H}$ 的绝对值在任一惯性系中彼此相等，那么它们在任何其他惯性系中也相等.

下面的不等式也显然成立.如果在任一参考系中 $E>H$（或 $H>E$）那么在每个其他参考系中也有 $E>H$（或
$H>E$）.如果在任一参考系中矢量 $\boldsymbol{E}$ 和 $\boldsymbol{H}$
成锐角（或钝角），那么它们在每个其他参考系中也成锐角（或钝角）.

借助洛伦兹变换，我们总可以给予 $\boldsymbol{E}$ 和 $\boldsymbol{H}$
任何任意的值，唯一的附加条件是 $E^2-H^2$ 和 $\boldsymbol{E}\cdot\boldsymbol{H}$
具有固定值.特别是，我们总可以找到一个惯性系，其中电场和磁场在给定点彼此平行.在这个参考系中
$\boldsymbol{E}\cdot\boldsymbol{H}=EH$，并且从下面两个方程
\begin{equation*}
  E^2-H^2=E_0^2-H_0^2,\qquad EH=\boldsymbol{E}_0\cdot\boldsymbol{H}_0
\end{equation*}
我们可以求出 $E$ 和 $H$ 在这个参考系中的值（$\boldsymbol{E}_0$ 和 $\boldsymbol{H}_0$
是在原来参考系中的电场和磁场）.

两个不变量都为零的情形除外.在这种情形下，$\boldsymbol{E}$ 和 $\boldsymbol{H}$
在所有参考系中都相等并且相互垂直.

如果 $\boldsymbol{E}\cdot\boldsymbol{H}=0$，我们总可以找到一个参考系，其中 $\boldsymbol{E}=0$ 或者
$\boldsymbol{H}=0$（依 $E^2-H^2<$ 或 $>0$ 而定），即纯磁场或纯电场.反之，如果在任一参考系中
$\boldsymbol{E}=0$ 或者 $\boldsymbol{H}=0$，则它们在每个其他参考系中都相互垂直，这与上节末的论述一致.

我们还将用另外的方法来求解反对称四维张量的不变量.特别是，从这个方法我们将看到，(25.3) 和 (25.4)
实际上是仅有的两个独立的不变量，同时我们将阐明，在应用于这样一种四维张量时，洛伦兹变换所具有的富有教益的数学性质.

我们来考虑复矢量
\begin{equation}
  \boldsymbol{F}=\boldsymbol{E}+\mathrm{i}\boldsymbol{H}.
\end{equation}

用公式 (24.2)—(24.3)，容易看出，对于这个矢量的洛伦兹变换（沿 $x$ 轴）具有形式
\begin{equation}
\begin{aligned}
  &F_x=F_x',\qquad
  F_y=F_y'\cosh\varphi-\mathrm{i}F_z'\sinh\varphi
  =F_y'\cos\mathrm{i}\varphi-F_z'\sin\mathrm{i}\varphi,\\
  &F_z=F_z'\cos\mathrm{i}\varphi+F_y'\sin\mathrm{i}\varphi,\qquad
  \tanh\varphi=\frac{V}{c}.
\end{aligned}
\end{equation}

我们看到，对于矢量 $\boldsymbol{F}$ 来说，在四维空间 $xt$
平面内的转动（就是这个洛伦兹变换），等价于在三维空间 $yz$
平面内转动一个虚角.四维空间内所有可能转动的集合（也包括绕 $x$，$y$ 和 $z$
轴的简单转动），等价于三维空间内转动一个复角的所有可能转动（这里四维空间内的6个转动角，相应于三维空间中的3个复转动角）.

矢量在转动下的唯一不变量是它的平方：
$\boldsymbol{F}^2=\boldsymbol{E}^2-\boldsymbol{H}^2+2\mathrm{i}\,\boldsymbol{E}\cdot\boldsymbol{H}$.因此，实量
$E^2-H^2$ 和 $\boldsymbol{E}\cdot\boldsymbol{H}$ 是张量 $F_{ik}$ 仅有的两个独立的不变量.

如果 $\boldsymbol{F}^2\neq0$，矢量 $\boldsymbol{F}$ 可以写成
$\boldsymbol{F}=a\boldsymbol{n}$，式中 $\boldsymbol{n}$ 是复单位矢量 $(n^2=1)$.通过适当的复转动，我们可以让
$\boldsymbol{n}$ 指向一条坐标轴；显然 $\boldsymbol{n}$ 变为实并决定了两个矢量 $\boldsymbol{E}$ 和
$\boldsymbol{H}$ 的方向：$\boldsymbol{F}=(E+\mathrm{i}H)\boldsymbol{n}$.换言之，我们得到
$\boldsymbol{E}$ 和 $\boldsymbol{H}$ 彼此平行的结果.

\begin{xiti}

\noindent{\heiti 习\ 题}\quad 有一个参考系，在其中，电场与磁场是相互平行的，试求这个参考系的速度.

解：满足所提条件的参考系 $K'$ 有无穷多.假如我们找到一个这样的参考系，那么，相对这个参考系的运动速度沿着
$\boldsymbol{E}$ 及 $\boldsymbol{H}$
的共同方向之任何其他参考系，也都具有同样的特性.因此，我们只需找到这些参考系中的一个就够了，这个参考系有与两个场都垂直的速度.选择速度的方向为
$x$ 轴的方向，利用这个事实，即在 $K'$ 中 $E_x'=H_x'=0$，
$E_y'H_z'-E_z'H_y'=0$，我们借助于公式（24.2）及（24.3）得到参考系 $K'$ 相对于原来参考系的速度
$\boldsymbol{V}$ 的方程如下：
\begin{equation*}
  \frac{\dfrac{V}{c}}{1+\dfrac{V^2}{c^2}}
  =\frac{\boldsymbol{E}\times\boldsymbol{H}}{E^2+H^2}
\end{equation*}
（在二次方程的两个根中自然应该选择 $V<c$ 的那一个根）.

\end{xiti}
